\subsection{注入延迟与单故障运行}
\label{sec:eval-faults}
在 \texttt{9879f64} 上，四种条件各运行三次，顺序交替。H 使用注入代理但不增加延迟；L 对成员 RPC 与委员会 P2P 流量每方向增加 25\,ms。L-M 与 L-C 增加该延迟，并分别在 15 至 45\,s 暂停成员~3 或委员会验证者~3。每次运行发送 100 笔预热支付、1,200 笔 20/s 的普通支付，以及两条在故障窗口内启动的 10 跳链。实际的健康检查 RPC RTT 从 1.17 升至 52.20\,ms；RPC 与 P2P 的平均保持时间均约为每方向 25.2\,ms。补充材料规定了 FIFO 注入、固定资源和逐次运行的结果。

全部 15,840 笔支付均闭合，委员会状态相等，发件箱为空，公开缺口为零。无需赔付。表~\ref{tab:eval-faults} 区分了接收、公开观测以及对全部四名成员的观测。暂停一名成员时，新的 QC 使用其余三名成员，公开完成继续进行，而缺席成员的观测造成 630--641 个未完成任务的峰值。暂停一名验证者时，接收 P95 保持在 108--111\,ms，每个故障窗口内有 400--510 笔支付达到公开观测，但公开 P95 升至 9.5--12.6\,s。保留的提交证书验证了每次运行中暂停验证者的五次第 0 轮提议者机会；每次都在第~1 轮提交，且不含该验证者的签名。两条链均在各自故障窗口内完成认证，且每次运行都在其最后一笔普通发送之后 1.25\,s 内排空。这些是单机注入延迟的结果，与 fast/wait 对比相互独立。

\begin{table}[t]
\caption{注入延迟与单故障：每次运行全部 1,200 笔普通支付的三个运行级统计量的中位数；预热与链不计入。公开与成员列为发送到观测的 P95；成员列等待全部四名成员。}
\label{tab:eval-faults}
\centering\small
\begin{tabular}{@{}lrrr@{}}
\toprule
条件 & 接收 P50/P95 (ms) & 公开 (s) & 成员 (s)\\
\midrule
H   & 41.56/83.20  & 1.28 & 1.36\\
L   & 81.91/118.22 & 1.32 & 1.41\\
L-M & 77.71/112.85 & 1.30 & 29.00\\
L-C & 77.41/109.32 & 10.02 & 10.04\\
\bottomrule
\end{tabular}
\end{table}
